% ==============================================================
\section{Proof of Deduplication Validity}\label{app:dedup_proof}
% ==============================================================

\begin{proof}[Proof of Claim~\ref{prop:dedup}]
Under the independent occurrence assumption,
the expected support of $P^*$ under the influence of event $e$ is $(p_{\text{base}} + \beta)^l$,
and $p_{\text{base}}^l$ outside the event period.
Setting $r = (p_{\text{base}} + \beta) / p_{\text{base}} > 1$,
we have $\text{magnitude}(P^*) \propto p_{\text{base}}^l (r^l - 1)$.
Secondary pattern $P'$ has $s$ boosted items and $l - s$ unboosted items, so
$\text{magnitude}(P') \propto p_{\text{base}}^l (r^s - 1)$.
Since $r > 1$ and $l > s$, we have $r^l > r^s$,
thus $\text{magnitude}(P^*) > \text{magnitude}(P')$ holds for any pattern length $l$.
Note that this result depends on the independent occurrence assumption;
when items have strong correlations, the ordering may reverse.
\end{proof}

% ==============================================================
\section{Synthetic Data Generation Details}\label{app:synthetic}
% ==============================================================

\subsection{Default Conditions}

The base vocabulary contains 200 items ($|\mathcal{I}| = 200$),
each with occurrence probability $p_{\text{base}} = 0.03$,
and $N = 100{,}000$ transactions.
The expected number of items per transaction is $200 \times 0.03 = 6.0$.

\begin{itemize}
\item \textbf{Type A} (boost, 5 patterns, mixed lengths):
    $L=2$ vocab-in patterns ($(5,15)$, $(25,35)$, $(45,55)$)
    have baseline support at $p_{\text{base}}$ (in-vocabulary).
    The $L=3$ pattern ($(201,202,203)$) and $L=4$ pattern ($(205,206,207,208)$)
    use out-of-vocabulary item IDs ($>200$) with zero background probability.
    Each pattern's event duration is 6{,}000 transactions (6\% of total),
    equally spaced on the time axis. During the event period,
    each item is additionally inserted with probability $\beta$.
\item \textbf{Type B} (event-unrelated, 2 patterns):
    Patterns $(65,75)$ and $(85,95)$ receive a boost of $\beta$
    in a random 6{,}000-transaction interval, but do not correspond to any event.
\item \textbf{Type C} (decoy events, 2):
    Events unrelated to any pattern variation (duration 6{,}000), randomly placed.
\end{itemize}

\subsection{Structural Conditions}

\begin{itemize}
\item \textbf{OVERLAP}: Events E1$=[16{,}000,28{,}000]$ and E2$=[24{,}000,36{,}000]$
    overlap in $[24{,}000,28{,}000]$. Tests temporal discrimination ability.
\item \textbf{CONFOUND}: The active interval of Type~B patterns is
    intentionally placed at event time windows (confounding condition).
\item \textbf{DENSE}: Doubled to Type~A 6 patterns + Type~B 4 patterns +
    4 decoys (high-density condition).
\item \textbf{SHORT}: Event duration shortened to 1{,}600
    (approximately 27\% of the default 6{,}000; short event condition).
\end{itemize}

\subsection{EX2 Ablation Scenarios}

\begin{itemize}
\item \textbf{Scenario A} (prox required):
    Pattern $(5,15)$ has two dense intervals---
    one near the event at $[15{,}000,25{,}000]$ (causal),
    and one at $[60{,}000,68{,}000]$ (accidental).
    Event interval: $[15{,}000,25{,}000]$.
\item \textbf{Scenario B} (mag required):
    Pattern $(5,15)$ is affected by two events---
    E1 with boost $\beta = 0.5$ (strong) and E2 with $\beta = 0.1$ (weak).
\end{itemize}

\subsection{Zipf Distribution Parameters}

Occurrence probability of item $k$ ($k = 1, \ldots, 200$):
$p(k) = \min(C / k^{\alpha_z},\; 0.10)$.
$C$ is set so that the median item ($k=100$) has probability 0.03,
with the maximum probability capped at 0.10.

% ==============================================================
\section{Dunnhumby Dataset Details}\label{app:dunnhumby}
% ==============================================================

An overview of the Dunnhumby ``The Complete Journey'' dataset~\cite{dunnhumby2014complete}
is given in Table~\ref{tab:dunnhumby_stats}.

\begin{table}[t]
\caption{Dunnhumby Dataset Statistics}
\label{tab:dunnhumby_stats}
\centering
\begin{tabular}{lr}
\toprule
Metric & Value \\
\midrule
Households & 2{,}500 \\
Transactions & 276{,}484 \\
Unique products & 92{,}339 \\
Mean items per transaction & 9.39 \\
Observation period & 711 days \\
Campaigns & 30 \\
\bottomrule
\end{tabular}
\end{table}

The 30 campaigns are classified into three types.
Based on the problem definition (Definition~\ref{def:attribution})
targeting population-level support variation,
this experiment uses only the 5 TypeA campaigns as external events
(see \S\ref{sec:exp_ex4}):
\begin{itemize}
\item \textbf{TypeA} (5 campaigns: C8, C13, C18, C26, C30): Large-scale promotions.
    3{,}700--38{,}200 coupon items, 40--56 day duration.
\item \textbf{TypeB} (18 campaigns): Targeted campaigns.
    200--2{,}000 coupon items, 32--61 day duration.
\item \textbf{TypeC} (7 campaigns): Niche campaigns.
    200--1{,}500 coupon items, 32--162 day duration.
\end{itemize}

Pipeline settings differ from Table~\ref{tab:default_params}:
$W = 300$ and $\theta = 5$.
$W$ was chosen as an appropriate temporal resolution
for 276K transactions.

% ==============================================================
\section{Null FDR Verification Details (100 Seeds)}\label{app:null_fdr}
% ==============================================================

We verified FDR control by the BH procedure under null conditions over 100 seeds.

\paragraph{Settings}
No planted signals (0 Type~A patterns),
5 randomly placed decoy events (each lasting 6,000 transactions),
4 Type~B dense patterns (temporally unrelated to any event),
$N = 100{,}000$, $W = 1{,}000$, $\theta = 100$, $B = 5{,}000$, $\alpha = 0.10$.

\paragraph{Results}

\begin{center}
\begin{tabular}{lr}
\toprule
Metric & Value \\
\midrule
Total seeds & 100 \\
Seeds with 0 false positives & 93/100 \\
Seeds with 1 false positive & 7/100 \\
Seeds with $\geq 2$ false positives & 0/100 \\
Mean false positive attributions & 0.07 \\
Mean per-seed FDR & 0.070 \\
FDR $\leq \alpha = 0.10$ & \textbf{YES} \\
\bottomrule
\end{tabular}
\end{center}

Mean per-seed FDR = 0.070 $\leq$ $\alpha = 0.10$,
confirming that the BH procedure controls FDR under null conditions over 100 seeds.
The 7 false positives (7\%) fall within the expected false discovery rate
$\alpha = 0.10$ of the BH procedure and are explained by stochastic variation.

% ==============================================================
\section{Robustness to Inter-Item Correlation}\label{app:item_correlation}
% ==============================================================

In real data, strong co-occurrence correlations (e.g., beer and snacks) exist,
which may violate the independent occurrence assumption of
Proposition~2 (validity of deduplication).
We performed 20-seed evaluation under both Zipf distribution and inter-item correlation conditions.

\paragraph{Settings}
\textbf{zipf\_only}: Zipf-1.0 distribution (independent occurrence).
\textbf{zipf\_corr}: Zipf-1.0 + correlation probability 0.5--0.7 assigned to 5 item pairs.
$N=100{,}000$, $W=1{,}000$, $\theta=100$, 20 seeds.

\paragraph{Results}
\begin{center}
\begin{tabular}{lccccc}
\toprule
Condition & Prec & Rec & F1 & 95\% CI & FAR \\
\midrule
zipf\_only & 0.74 & 0.92 & 0.81 & [0.70, 0.92] & 0.25 \\
zipf\_corr & 0.41 & 0.85 & 0.55 & [0.45, 0.65] & 0.20 \\
\bottomrule
\end{tabular}
\end{center}

F1 = 0.81 for zipf\_only, dropping to F1 = 0.55 for zipf\_corr.
The primary cause of accuracy degradation is that additional dense patterns
generated by correlated item pairs increase false positive attributions.
This supports the possibility that the independent occurrence assumption
of Proposition~2 is relaxed in real data;
FAR = 0.20 under strong inter-item correlation remains acceptable.

% ==============================================================
\section{Cross-Signal Effects in OVERLAP Condition}\label{app:overlap_analysis}
% ==============================================================

We explain why F1 variance is high (CI width = 0.44) in the OVERLAP condition
(\S\ref{sec:exp_ex1}).
In the overlap interval $[24{,}000, 28{,}000]$ where
E1 $= [16{,}000, 28{,}000]$ and E2 $= [24{,}000, 36{,}000]$ overlap,
the item sets of E1 and E2 are \emph{simultaneously boosted within the same window}.
With $\beta = 0.3$, the expected support of cross-patterns
(combinations of E1 and E2 items) is
$0.3 \times 0.3 \times W = 90 \approx \theta$,
and statistical variation causes short dense intervals that exceed $\theta$.
Deduplication (Step~5) eliminates most of these by merging them
with strong planted patterns, but in some seeds cross-patterns remain as false positives,
causing variance in accuracy.

% ==============================================================
\section{Parameter Selection Rationale}\label{app:params}
% ==============================================================

\begin{table}[h]
\caption{Default Parameters}
\label{tab:default_params}
\centering
\begin{tabular}{ll}
\toprule
Parameter & Value \\
\midrule
Window size $W$ & 1{,}000 \\
Min. support $\theta$ & 100 \\
Attribution target patterns & $|P| \geq 2$ \\
Max. pattern length $L$ & unconstrained\footnotemark[1] \\
Change point detection & threshold crossing \\
$\sigma$ & $W$ (= 1{,}000) \\
Permutations $B$ & 5{,}000 \\
Significance level $\alpha$ & 0.10 \\
Correction method & BH \\
Magnitude normalization & square root (Eq.~(\ref{eq:mag_sqrt})) \\
Deduplication & enabled \\
\bottomrule
\end{tabular}
\end{table}

\footnotetext[1]{Due to Apriori's anti-monotonicity, the search naturally terminates
when no patterns satisfying the support threshold exist.}

Rationale for key parameter choices:

\begin{itemize}
\item \textbf{Window size $W = 1{,}000$}:
    Corresponds to 1\% temporal resolution for $N = 100{,}000$ transactions.
    Too small a $W$ makes support estimation unstable;
    too large reduces temporal precision of change point detection.
\item \textbf{Min. support $\theta = 100$}:
    Sufficiently large relative to the expected baseline support
    $W \cdot p_{\text{base}}^2 = 1{,}000 \times 0.03^2 = 0.9$,
    excluding chance co-occurrence while capturing the support increase during boost.
    $\theta = W / 10 = 100$ suppresses the number of false positive patterns
    under Zipf distribution.
\item \textbf{Decay parameter $\sigma = W$}:
    Proximity decays to $e^{-1} \approx 0.37$ at a distance of one window width.
    Attribution score decreases rapidly when the change point-to-event boundary
    distance exceeds $W$.
\item \textbf{Permutations $B = 5{,}000$}:
    $p$-value resolution $1/(B+1) \approx 0.0002$.
    Guarantees sufficient $p$-value precision even when testing hundreds of hypotheses
    simultaneously at $\alpha = 0.10$ with BH correction.
\item \textbf{Significance level $\alpha = 0.10$}:
    Standard FDR control level.
    Compared to $\alpha = 0.05$, this setting prioritizes detection power,
    tolerating a false discovery rate below 10\%.
\item \textbf{Magnitude normalization}:
    Square root normalization (Eq.~(\ref{eq:mag_sqrt})) is adopted.
    Absolute magnitude has a bias where patterns with higher baseline inflate scores;
    full normalization ($\div s^{-}$) over-amplifies scores of low-baseline patterns.
    Square root normalization lies between these extremes and achieves
    stable accuracy for real-world data with Zipf-like frequency distributions.
\end{itemize}
